\PassOptionsToPackage{unicode=true}{hyperref} % options for packages loaded elsewhere
\PassOptionsToPackage{hyphens}{url}
%
\documentclass[]{book}
\usepackage{lmodern}
\usepackage{amssymb,amsmath}
\usepackage{ifxetex,ifluatex}
\usepackage{fixltx2e} % provides \textsubscript
\ifnum 0\ifxetex 1\fi\ifluatex 1\fi=0 % if pdftex
  \usepackage[T1]{fontenc}
  \usepackage[utf8]{inputenc}
  \usepackage{textcomp} % provides euro and other symbols
\else % if luatex or xelatex
  \usepackage{unicode-math}
  \defaultfontfeatures{Ligatures=TeX,Scale=MatchLowercase}
\fi
% use upquote if available, for straight quotes in verbatim environments
\IfFileExists{upquote.sty}{\usepackage{upquote}}{}
% use microtype if available
\IfFileExists{microtype.sty}{%
\usepackage[]{microtype}
\UseMicrotypeSet[protrusion]{basicmath} % disable protrusion for tt fonts
}{}
\IfFileExists{parskip.sty}{%
\usepackage{parskip}
}{% else
\setlength{\parindent}{0pt}
\setlength{\parskip}{6pt plus 2pt minus 1pt}
}
\usepackage{hyperref}
\hypersetup{
            pdftitle={Instalación de R y RStudio},
            pdfauthor={Eva María Mazcuñán Navarro},
            pdfborder={0 0 0},
            breaklinks=true}
\urlstyle{same}  % don't use monospace font for urls
\usepackage{longtable,booktabs}
% Fix footnotes in tables (requires footnote package)
\IfFileExists{footnote.sty}{\usepackage{footnote}\makesavenoteenv{longtable}}{}
\usepackage{graphicx,grffile}
\makeatletter
\def\maxwidth{\ifdim\Gin@nat@width>\linewidth\linewidth\else\Gin@nat@width\fi}
\def\maxheight{\ifdim\Gin@nat@height>\textheight\textheight\else\Gin@nat@height\fi}
\makeatother
% Scale images if necessary, so that they will not overflow the page
% margins by default, and it is still possible to overwrite the defaults
% using explicit options in \includegraphics[width, height, ...]{}
\setkeys{Gin}{width=\maxwidth,height=\maxheight,keepaspectratio}
\setlength{\emergencystretch}{3em}  % prevent overfull lines
\providecommand{\tightlist}{%
  \setlength{\itemsep}{0pt}\setlength{\parskip}{0pt}}
\setcounter{secnumdepth}{5}
% Redefines (sub)paragraphs to behave more like sections
\ifx\paragraph\undefined\else
\let\oldparagraph\paragraph
\renewcommand{\paragraph}[1]{\oldparagraph{#1}\mbox{}}
\fi
\ifx\subparagraph\undefined\else
\let\oldsubparagraph\subparagraph
\renewcommand{\subparagraph}[1]{\oldsubparagraph{#1}\mbox{}}
\fi

% set default figure placement to htbp
\makeatletter
\def\fps@figure{htbp}
\makeatother

\usepackage{booktabs}
\usepackage{etoolbox}
\makeatletter
\providecommand{\subtitle}[1]{% add subtitle to \maketitle
  \apptocmd{\@title}{\par {\large #1 \par}}{}{}
}
\makeatother
\usepackage[]{natbib}
\bibliographystyle{apalike}

\title{Instalación de R y RStudio}
\providecommand{\subtitle}[1]{}
\subtitle{Métodos Numéricos y Estadísticos}
\author{Eva María Mazcuñán Navarro}
\date{}

\begin{document}
\maketitle

{
\setcounter{tocdepth}{1}
\tableofcontents
}
\hypertarget{section}{%
\chapter*{}\label{section}}
\addcontentsline{toc}{chapter}{}

Utilizaremos el lenguaje \textbf{R}, vía la IDE \textbf{RStudio}, para resolver las tareas que se plantearán en la asignatura.

En este libro encontrarás las instrucciones para instalar R y RStudio en tu equipo.

Utiliza la tabla de contenidos del panel izquierdo para elegir el sistema operativo donde quieres realizar la instalación, y sigue los pasos que se describen allí.

\hypertarget{linux}{%
\chapter{Linux}\label{linux}}

\begin{center}\includegraphics[width=0.15\linewidth]{images/os/tux-flat} \end{center}

\hypertarget{instalaciuxf3n-de-r}{%
\section{Instalación de R}\label{instalaciuxf3n-de-r}}

En este apartado se describen las instrucciones para instalar R en Ubuntu. Para otras distribuciones de Linux consulta \href{https://ftp.cixug.es/CRAN/bin/linux/}{este link}.

\begin{center}\includegraphics[width=0.15\linewidth]{images/os/ubuntu} \end{center}

Para instalar R en Ubuntu instala los paquetes \texttt{r-base} y \texttt{r-base-dev}, ejecutando desde la terminal

\begin{verbatim}
sudo apt-get update
sudo apt install r-base r-base-dev 
\end{verbatim}

También puedes hacer la instalación desde \textbf{Ubuntu Software}.

\hypertarget{instalaciuxf3n-de-rstudio-versiuxf3n-desktop-free}{%
\section{Instalación de RStudio (versión Desktop Free)}\label{instalaciuxf3n-de-rstudio-versiuxf3n-desktop-free}}

Una vez instalado R, sigue los siguientes pasos para instalar RStudio.

\textbf{PASO 1: }
Accede a la página de descargas de RStudio pinchando \href{https://rstudio.com/products/rstudio/download/\#download}{aquí}.

En el apartado \textbf{All Installers} descarga el instalador más reciente para tu versión (para Ubuntu 20, descarga Ubuntu 18/Debian 10).

\textbf{PASO 2: }
Cuando se complete la descarga del paso anterior, haz doble click en el archivo .deb descargado. Se abrirá con el gestor de software de Ubuntu. Presiona el botón `Instalar' para terminar.

\hypertarget{windows}{%
\chapter{Windows}\label{windows}}

\begin{center}\includegraphics[width=0.15\linewidth]{images/os/windows} \end{center}

A continuación se proporcionan los links para la descarga de los instaladores de R y RStudio en Windows. En ambos casos, se descargará un archivo ejecutable que gestionará la instalación de R / RStudio en tu equipo. Para ejecutarlo, se recomienda hacer click derecho en el archivo y seleccionar `Ejecutar como administrador'. Puedes conservar todas las opciones de instalación que aparecen por defecto.

\hypertarget{instalaciuxf3n-de-r-1}{%
\section{Instalación de R}\label{instalaciuxf3n-de-r-1}}

Accede a la página de descargas de R para Windows pinchando \href{https://ftp.cixug.es/CRAN/bin/windows/base/}{aquí}, y haz click el link de título `Download R 4.0.4 for Windows', que encontrarás al principio de la página.

\hypertarget{instalaciuxf3n-de-rstudio-desktop-free}{%
\section{Instalación de RStudio (Desktop Free)}\label{instalaciuxf3n-de-rstudio-desktop-free}}

Accede a la página de descargas de RStudio pinchando \href{https://rstudio.com/products/rstudio/download/\#download}{aquí}, y, en la tabla del apartado \textbf{All Installers}, haz click en el link de descarga correspondiente a Windows 10/8/7.

\hypertarget{mac-os-x}{%
\chapter{Mac OS X}\label{mac-os-x}}

\begin{center}\includegraphics[width=0.15\linewidth]{images/os/apple} \end{center}

\hypertarget{instalaciuxf3n-de-r-2}{%
\section{Instalación de R}\label{instalaciuxf3n-de-r-2}}

Accede a la página de descargas de R para Mac OS X pinchando \href{https://ftp.cixug.es/CRAN/bin/macosx/}{aquí} y descarga el instalador correspondiente a la versión de tu equipo.

Solo tienes que ejecutar el archivo pgk que se descargará para instalar R en tu equipo. Puedes conservar todas las opciones de instalación por defecto.

\hypertarget{instalaciuxf3n-de-rstudio}{%
\section{Instalación de RStudio}\label{instalaciuxf3n-de-rstudio}}

Accede a la página de descargas de RStudio pinchando \href{https://rstudio.com/products/rstudio/download/\#download}{aquí}, y, en la tabla del apartado \textbf{All Installers}, haz click en el link de descarga correspondiente a macOS 10.13+.

Se descargará un archivo dmg. Para completar la instalación de RStudio en tu equipo solo tienes que arrastar el icono de la aplicación a tu carpeta Aplicaciones.

\bibliography{book.bib,packages.bib}

\end{document}
